\section{Conclusion}
\label{sec:conclusion}

A multi-fidelity aerodynamic shape optimisation framework has been developed and demonstrated on the Ahmed body ($25^\circ$ slant baseline, $\mathrm{Re}=2.78\times10^6$). The pipeline combines a two-phase RANS design-of-experiments (80 L2 evaluations), 20 Bayesian optimisation iterations on the L2 surrogate, and an additive Kennedy \& O'Hagan co-Kriging correction trained on 8 L3 anchor simulations. The verified multi-fidelity optimum is confirmed by a dedicated L3 CFD solve.

\textbf{Key findings:}

\begin{itemize}
    \item \textbf{Single-fidelity optimum (L2 RANS):} $\alpha_\mathrm{slant}=32.5^\circ$, $\alpha_\mathrm{diff}=8.9^\circ$, $C_d=0.319$, $C_l=-0.327$, $f_1=0.210$. Sobol sensitivity analysis eliminated nose radius and fixed ride height; slant and diffuser angles together account for $>97\%$ of objective variance.

    \item \textbf{Fidelity gap:} An L3 validation run at the L2 optimum revealed $\Delta C_l=+0.160$ (L3 predicts substantially less downforce), inflating $f_1$ to $0.248$ — an 18\% degradation. The gap is spatially non-uniform: $\delta_{C_l}$ ranges from $-0.114$ to $+0.167$ across the 8 anchor locations, demonstrating that a constant-offset correction would mis-locate the HF minimum.

    \item \textbf{Multi-fidelity optimum (verified):} $\alpha_\mathrm{slant}=38.9^\circ$, $\alpha_\mathrm{diff}=10.5^\circ$, $C_d=0.292$, $C_l=-0.285$, $f_1=0.198$ (L3 CFD, $\sigma_{C_d}=0.001$). The surrogate predicted $f_1=0.196$; the L3 verification returned $f_1=0.198$ — a $0.6\%$ error within the $\pm2\sigma_\mathrm{HF}$ band. This represents a 5.7\% improvement over the L2 optimum and a 20.2\% improvement over L3 evaluated at the old optimum location.

    \item \textbf{Optimum robustness:} A $\lambda$-sensitivity sweep shows the MF optimum location $(38.8$--$39.0^\circ$, $10.4$--$10.5^\circ$) is stable for any $\lambda>0$ (any positive downforce weighting). Only the pure drag-minimisation case ($\lambda=0$) relocates the optimum to the low-slant regime. Basin-mapping L3 stencil solves confirm a smooth gradient from the critical-angle region toward the optimum: $f_1=0.215$ at the cliff point ($35.7^\circ/9.3^\circ$), transitioning smoothly to $f_1=0.198$ at the optimum.

    \item \textbf{BO iter~1 artefact resolved:} The anomalous training point ($f_\mathrm{cfd}=0.642$ at $35.7^\circ/9.3^\circ$) was a pseudo-time limit-cycle oscillation on the coarse L2 mesh ($\sigma_{C_d}=0.023$). An L3 re-solve at the same geometry returns $f_1=0.215$ with $\sigma_{C_d}=0.002$, confirming the value is physically reasonable. The L2 artefact inflated GP uncertainty in the transition region, biasing BO acquisition away from the high-slant basin where the true optimum resides.

    \item \textbf{Scope of the correction:} The co-Kriging correction spans mesh fidelity only (L2 $\to$ L3 steady RANS). Turbulence-model error from the steady kOmegaSST closure — including incorrect representation of the bistable critical-angle regime — is common to both fidelity levels and not corrected. The present result is the optimum of the L3 RANS landscape; URANS or DES anchors represent the natural next fidelity axis.
\end{itemize}

\textbf{Compute efficiency:} 100 L2 RANS cases (${\sim}$200~s each, 8 cores) + 9 L3 RANS cases (${\sim}$500~s each, 8 cores) = ${\sim}$60 CPU-hours total for the full optimisation campaign. The co-Kriging framework delivers a verified, fidelity-corrected optimum at a fraction of the cost of a full L3 design sweep.
